% Generated draft: requires mathematical and source review.
\documentclass[11pt]{article}
\usepackage[margin=25mm]{geometry}
\usepackage{fontspec}
\setmainfont{Latin Modern Roman}
\usepackage{amsmath,amssymb,mathtools,unicode-math}
\setmathfont{Latin Modern Math}
\usepackage{enumitem,needspace,xurl,hyperref}
\setlength{\parindent}{0pt}
\setlength{\parskip}{5pt}
\setlength{\emergencystretch}{3em}
\newcommand{\N}{\mathbb{N}}
\newcommand{\Z}{\mathbb{Z}}
\newcommand{\Q}{\mathbb{Q}}
\newcommand{\R}{\mathbb{R}}
\newcommand{\C}{\mathbb{C}}
\newcommand{\F}{\mathbb{F}}
\newcommand{\A}{\mathbb{A}}
\newcommand{\PP}{\mathbb{P}}
\newcommand{\E}{\mathbb{E}}
\newcommand{\eps}{\varepsilon}
\newcommand{\dd}{\,\mathrm{d}}
\begin{document}
\section*{UnsolvedMath ID: 701639}
\textbf{Source problem code:} catalogue entry
\subsection{Definitions and mathematical statement}
Let \(p\) be a prime, let \(G\) be a finite group, and let \(A\) be an abelian Sylow \(p\)-subgroup of \(G\). A subgroup \(B\leq A\) is elementary abelian if it is abelian and every nonidentity element has order \(p\). Here strong closure of \(B\) in \(A\) with respect to \(G\) means
\[
\forall b\in B\;\forall g\in G,\qquad g^{-1}bg\in A\ \Longrightarrow\ g^{-1}bg\in B.
\]
For an abelian \(p\)-group \(P\), define
\[
\Omega_1(P)=\langle x\in P:x^p=1\rangle=\{x\in P:x^p=1\}.
\]
\textbf{Requested statement and proof requirement.} For every prime \(p\), every finite group \(G\), every abelian Sylow \(p\)-subgroup \(A\) of \(G\), and every strongly closed elementary abelian subgroup \(B\leq A\), prove, without invoking the Classification Theorem for Finite Simple Groups (CFSG), that
\[
\exists N\trianglelefteq G\qquad B=\Omega_1(A\cap N).
\]
The CFSG-free requirement concerns the proof method and is essential to the supplied problem.

The source states that the case \(p=2\) is a corollary of a theorem of Goldschmidt, and that Flores and Foote proved the assertion for odd \(p\), using CFSG. These are supplied literature assertions, not independently verified here.

The source also singles out a CFSG-free proof under the additional hypotheses
\[
p=3,\qquad \operatorname{rank}_3(A)=3,
\]
where \(\operatorname{rank}_3(A)=\dim_{\F_3}\Omega_1(A)\), equivalently the largest dimension of an elementary abelian subgroup of \(A\). It describes this special case as already quite interesting because it arises in Aschbacher's treatment of the \(e(G)=3\) problem and would provide an alternative to part of his argument. For that application, the source permits assuming that all proper simple sections of \(G\) are known. A simple section means a simple quotient \(H/K\) with \(K\trianglelefteq H\leq G\). The source does not specify the meaning of \(e(G)\), which sections are intended by ``proper,'' or the precise information meant by ``known''; this application-specific allowance is therefore retained as an unresolved qualification, not added as a formal hypothesis to the general statement.
\subsection{Short English statement}
Given an abelian Sylow \(p\)-subgroup \(A\) of a finite group \(G\) and a strongly closed elementary abelian subgroup \(B\leq A\), find a normal subgroup \(N\) whose elements of order dividing \(p\) in \(A\cap N\) are exactly \(B\). The requested proof must avoid CFSG. The source particularly highlights \(p=3\) with \(A\) of \(3\)-rank \(3\).
\subsection{Sources}
\begin{itemize}
\item Supplied local UnsolvedMath registry record, permanent ID \(701639\), exact source problem code \texttt{catalogue entry}, titled ``Kourovka Notebook Problem 19.91.'' The record attributes the problem to R. Solomon, proposed in 2018, and labels it \texttt{open}; this is a supplied status label, not an independently verified assessment. The record was created and updated on 2026-10-16.
\item \emph{The Kourovka Notebook (Unsolved Problems in Group Theory)}, original Problem 19.91, 19th issue (2018); the supplied background also identifies the 21st issue edition (2026), \texttt{arXiv:1401.0300v48}. The supplied collection description names E. I. Khukhro and V. D. Mazurov as editors. This source identification is supplied and has not been independently checked.
\item Flores and Foote, \emph{Adv. Math.} \textbf{222} (2009), 453--484. The source attributes the odd-prime result to this work and states that its proof relies on CFSG. No article title is supplied.
\item Goldschmidt: the source states that the \(p=2\) case follows from a theorem of his, but supplies no precise theorem reference or bibliographic details.
\item Aschbacher: the source mentions his treatment of the \(e(G)=3\) problem and the application-specific allowance concerning proper simple sections, but supplies no precise bibliographic reference.
\end{itemize}
\end{document}
